\chapter{Lesson 8: Verbs, Particles and Time Expressions}

\section{Japanese sentence structure}

Japanese is a \textbf{subject--object--verb} (SOV) language: the verb always comes last.
Particles after each noun tell us that noun's role, so the order of the noun phrases
before the verb is fairly free.

\begin{center}
\begin{tikzpicture}[b/.style={draw, rounded corners, minimum height=0.85cm, inner sep=4pt},
    p/.style={draw, circle, fill=orange!15, inner sep=2pt, minimum size=0.75cm}]
  \node[b, fill=blue!6] (t) at (0,0) {\ja{わたし}};
  \node[p, right=0.05cm of t] (wa) {\ja{は}};
  \node[b, fill=blue!6, right=0.25cm of wa] (ti) {\ja{あした}};
  \node[b, fill=blue!6, right=0.25cm of ti] (pl) {\ja{としょかん}};
  \node[p, right=0.05cm of pl] (de) {\ja{で}};
  \node[b, fill=blue!6, right=0.25cm of de] (o) {\ja{ほん}};
  \node[p, right=0.05cm of o] (wo) {\ja{を}};
  \node[b, fill=green!10, right=0.25cm of wo] (v) {\ja{よみます}};
  \node[font=\scriptsize, below=0.15cm of wa] {topic};
  \node[font=\scriptsize, below=0.15cm of ti] {time};
  \node[font=\scriptsize, below=0.15cm of de] {place of action};
  \node[font=\scriptsize, below=0.15cm of wo] {object};
  \node[font=\scriptsize, below=0.15cm of v] {verb (last)};
\end{tikzpicture}
\end{center}
\ja{わたしはあしたとしょかんでほんをよみます。}
\emph{Watashi wa ashita toshokan de hon o yomimasu.}
``I will read a book at the library tomorrow.''

\section{Verb groups}

Every verb has a \emph{dictionary form} (the form listed in a dictionary, ending in
-\emph{u}). Verbs fall into three groups by how they conjugate:
\begin{itemize}
  \item \textbf{\ja{る}-verbs} (ichidan, ``Group 2''): stem + \ja{る}; drop \ja{る} and add
        endings. \ja{たべる} $\to$ \ja{たべます}.
  \item \textbf{\ja{う}-verbs} (godan, ``Group 1''): the final \emph{u} changes to
        \emph{i} before \ja{ます}. \ja{のむ} $\to$ \ja{のみます}.
  \item \textbf{irregular verbs}: \ja{する} ``do'' and \ja{くる} ``come'' (and compounds
        like \ja{べんきょうする}).
\end{itemize}

\begin{center}
\begin{tikzpicture}[node distance=0.9cm and 1.2cm,
    q/.style={draw, diamond, aspect=2.4, fill=yellow!10, inner sep=1pt, font=\small, align=center},
    r/.style={draw, rounded corners, fill=blue!8, font=\small, align=center, minimum width=2.4cm, minimum height=0.8cm}]
  \node[q] (q1) {\ja{する} or \ja{くる}?};
  \node[r, right=of q1, fill=red!10] (irr) {irregular};
  \node[q, below=of q1] (q2) {ends in\\ -iru or -eru?};
  \node[r, right=of q2, fill=green!10] (u1) {\ja{う}-verb};
  \node[q, below=of q2] (q3) {on the list of\\ exceptions?};
  \node[r, right=of q3, fill=green!10] (u2) {\ja{う}-verb};
  \node[r, below=of q3] (ru) {\ja{る}-verb};
  \draw[-Stealth] (q1) -- node[above, font=\scriptsize] {yes} (irr);
  \draw[-Stealth] (q1) -- node[left, font=\scriptsize] {no} (q2);
  \draw[-Stealth] (q2) -- node[above, font=\scriptsize] {no} (u1);
  \draw[-Stealth] (q2) -- node[left, font=\scriptsize] {yes} (q3);
  \draw[-Stealth] (q3) -- node[above, font=\scriptsize] {yes} (u2);
  \draw[-Stealth] (q3) -- node[left, font=\scriptsize] {no} (ru);
\end{tikzpicture}
\end{center}

\begin{important}[Exceptions: look like \ja{る}-verbs but are \ja{う}-verbs]
\ja{かえる} kaeru (return), \ja{はいる} hairu (enter), \ja{はしる} hashiru (run),
\ja{しる} shiru (know), \ja{いる} iru (need --- not ``exist''), \ja{きる} kiru (cut).
Their \ja{ます}-forms: \ja{かえります}, \ja{はいります}, \ldots
\end{important}

\section{The \ja{ます}-form: present (and future) tense}

Japanese verbs have no separate future tense: the ``present'' form covers habitual
actions and future actions. With an action verb it never means ``is doing now'' (that is
\ja{〜ている}, Lesson 11).

\begin{center}
\begin{tabular}{lllll}
\toprule
Group & Dictionary & Stem & Affirmative & Negative \\
\midrule
\ja{る} & \ja{たべる} eat & \ja{たべ} & \ja{たべます} & \ja{たべません} \\
\ja{る} & \ja{みる} see, watch & \ja{み} & \ja{みます} & \ja{みません} \\
\ja{る} & \ja{おきる} get up & \ja{おき} & \ja{おきます} & \ja{おきません} \\
\ja{る} & \ja{ねる} sleep & \ja{ね} & \ja{ねます} & \ja{ねません} \\
\midrule
\ja{う} & \ja{いく} go & \ja{いき} & \ja{いきます} & \ja{いきません} \\
\ja{う} & \ja{のむ} drink & \ja{のみ} & \ja{のみます} & \ja{のみません} \\
\ja{う} & \ja{よむ} read & \ja{よみ} & \ja{よみます} & \ja{よみません} \\
\ja{う} & \ja{はなす} speak & \ja{はなし} & \ja{はなします} & \ja{はなしません} \\
\ja{う} & \ja{きく} listen & \ja{きき} & \ja{ききます} & \ja{ききません} \\
\ja{う} & \ja{かう} buy & \ja{かい} & \ja{かいます} & \ja{かいません} \\
\ja{う} & \ja{まつ} wait & \ja{まち} & \ja{まちます} & \ja{まちません} \\
\ja{う} & \ja{かえる} return & \ja{かえり} & \ja{かえります} & \ja{かえりません} \\
\midrule
irr. & \ja{する} do & \ja{し} & \ja{します} & \ja{しません} \\
irr. & \ja{くる} come & \ja{き} & \ja{きます} & \ja{きません} \\
\bottomrule
\end{tabular}
\end{center}

\begin{remark}
\ja{う}-verb endings shift along the row of the kana chart:
\ja{く}$\to$\ja{き}, \ja{む}$\to$\ja{み}, \ja{す}$\to$\ja{し}, \ja{つ}$\to$\ja{ち},
\ja{う}$\to$\ja{い}, \ja{る}$\to$\ja{り}. This is why knowing the chart by rows pays off.
\end{remark}

\section{Particles}

\subsection{\ja{を} (o): direct object}
\ja{コーヒーをのみます。} \emph{Kōhī o nomimasu.} (I) drink coffee.\\
\ja{テレビをみません。} \emph{Terebi o mimasen.} (I) don't watch TV.

\subsection{\ja{で} (de): place of action; means}
\ja{としょかんでべんきょうします。} \emph{Toshokan de benkyō shimasu.}
(I) study at the library.\\
\ja{にほんごではなします。} \emph{Nihongo de hanashimasu.} (We) speak in Japanese.\\
\ja{バスでいきます。} \emph{Basu de ikimasu.} (I) go by bus.

\subsection{\ja{に} (ni): goal; specific time}
\ja{だいがくにいきます。} \emph{Daigaku ni ikimasu.} (I) go to the university.\\
\ja{しちじにおきます。} \emph{Shichiji ni okimasu.} (I) get up at seven.\\
\ja{にちようびにかいものをします。} \emph{Nichiyōbi ni kaimono o shimasu.}
(I) go shopping on Sunday.

\subsection{\ja{へ} (e): direction}
\ja{へ} marks the direction of movement and is mostly interchangeable with goal
\ja{に} for \ja{いく}, \ja{くる}, \ja{かえる}: \ja{うちへかえります。} ``(I) go
home.'' It cannot replace \ja{に} for time.

\begin{important}[Time words with and without \ja{に}]
\begin{itemize}
  \item Clock times, days of the week, dates take \ja{に}: \ja{くじに}, \ja{どようびに}.
  \item Relative time words take \textbf{no} particle: \ja{きょう} today, \ja{あした}
        tomorrow, \ja{まいにち} every day, \ja{こんばん} tonight, \ja{いつ} when.
  \item \ja{〜ごろ} ``about'' may replace \ja{に}: \ja{じゅうにじごろねます。}
\end{itemize}
\end{important}

\section{Time expressions}
\begin{center}
\begin{tabular}{lll@{\qquad}lll}
\toprule
\ja{あさ} & asa & morning & \ja{きょう} & kyō & today \\
\ja{ひる} & hiru & noon & \ja{あした} & ashita & tomorrow \\
\ja{ばん} & ban & evening & \ja{こんばん} & konban & tonight \\
\ja{しゅうまつ} & shūmatsu & weekend & \ja{まいにち} & mainichi & every day \\
\ja{まいばん} & maiban & every night & \ja{まいあさ} & maiasa & every morning \\
\bottomrule
\end{tabular}
\end{center}

\subsection{Frequency adverbs}
\begin{center}
\begin{tabular}{lll}
\toprule
Adverb & Meaning & Verb form \\
\midrule
\ja{いつも} & always & affirmative \\
\ja{よく} & often & affirmative \\
\ja{ときどき} & sometimes & affirmative \\
\ja{あまり} & not much & \textbf{negative} \\
\ja{ぜんぜん} & not at all & \textbf{negative} \\
\bottomrule
\end{tabular}
\end{center}
\ja{ぜんぜんテレビをみません。} ``I don't watch TV at all.''

\subsection{Invitations: \ja{〜ませんか} and \ja{〜ましょう}}
\ja{いっしょにえいがをみませんか。} \emph{Issho ni eiga o mimasen ka.}
``Won't you see a movie with me?'' (negative question = invitation)\\
\ja{ええ、いいですね。みましょう。} \emph{Ee, ii desu ne. Mimashō.}
``Sounds good. Let's see it.''\\
To decline politely, trail off: \ja{きょうはちょっと\ldots} ``Today is a little\ldots''

\begin{aside}
Japanese speakers rarely say a flat ``no'' (\ja{いいえ}) to an invitation. The vague
\ja{ちょっと\ldots} (``it's a bit\ldots'') with an apologetic tone is understood as a
refusal; pressing for a reason is considered pushy.
\end{aside}

\section{Kanji: days of the week}

The days of the week are named after the sun, moon and five elements. These are the
first kanji we are responsible for writing.

\begin{center}
\begin{tabular}{llllll}
\toprule
Kanji & Meaning & On reading & Kun reading & Day & \\
\midrule
\ja{日} & sun, day & \ja{ニチ}, \ja{ジツ} & \ja{ひ}, \ja{び}, \ja{か} & \ja{日曜日} & nichiyōbi, Sunday \\
\ja{月} & moon, month & \ja{ゲツ}, \ja{ガツ} & \ja{つき} & \ja{月曜日} & getsuyōbi, Monday \\
\ja{火} & fire & \ja{カ} & \ja{ひ} & \ja{火曜日} & kayōbi, Tuesday \\
\ja{水} & water & \ja{スイ} & \ja{みず} & \ja{水曜日} & suiyōbi, Wednesday \\
\ja{木} & tree & \ja{モク} & \ja{き} & \ja{木曜日} & mokuyōbi, Thursday \\
\ja{金} & gold, money & \ja{キン} & \ja{かね} & \ja{金曜日} & kin'yōbi, Friday \\
\ja{土} & earth, soil & \ja{ド} & \ja{つち} & \ja{土曜日} & doyōbi, Saturday \\
\bottomrule
\end{tabular}
\end{center}

\begin{definition}[On and kun readings]
Most kanji have an \emph{on} reading (\ja{音読み}, borrowed from Chinese; written in
katakana in dictionaries) used mainly in compounds, and a \emph{kun} reading
(\ja{訓読み}, the native Japanese word; written in hiragana) used when the kanji stands
alone or with okurigana.
\end{definition}

\begin{center}
\begin{tikzpicture}[k/.style={draw, minimum size=1.3cm, font=\huge}]
  \node[k] (a) at (0,0) {\ja{木}};
  \node[k] (b) at (3,0) {\ja{林}};
  \node[k] (c) at (6,0) {\ja{森}};
  \draw[-Stealth, thick] (a) -- node[above, font=\small] {$\times 2$} (b);
  \draw[-Stealth, thick] (b) -- node[above, font=\small] {$\times 3$} (c);
  \node[font=\small] at (0,-1) {tree (\ja{き})};
  \node[font=\small] at (3,-1) {grove (\ja{はやし})};
  \node[font=\small] at (6,-1) {forest (\ja{もり})};
\end{tikzpicture}
\end{center}
(\ja{林} and \ja{森} are not on the quiz, but show how kanji can be built from components.)

\section*{Lesson 8 review}
\begin{itemize}
  \item SOV order; particles mark roles.
  \item \ja{る}-verbs: drop \ja{る} + \ja{ます}. \ja{う}-verbs: \emph{u}$\to$\emph{i} +
        \ja{ます}. Irregular: \ja{します}, \ja{きます}. Memorise the \ja{かえる}-type
        exceptions.
  \item Negative: \ja{〜ません}. Present form also covers the future.
  \item \ja{を} object, \ja{で} place of action/means, \ja{に} goal/time, \ja{へ}
        direction.
  \item No \ja{に} after \ja{きょう}, \ja{あした}, \ja{まいにち}.
  \item \ja{あまり}/\ja{ぜんぜん} + negative. Invitations: \ja{〜ませんか};
        \ja{〜ましょう}.
  \item Kanji: \ja{日月火水木金土}; on vs.\ kun readings.
\end{itemize}
